\subsection{The Role of Transparency and Accountability in Bias Reduction}
\label{sec:6.1.2}
Transparency and accountability address bias from two directions: transparency lets stakeholders see what is driving an AI system's outcomes, and accountability gives developers and users a reason to act on what transparency reveals rather than just document it.

\runin{Transparency in Research and Bias Auditing}
Section~\ref{sec:9.3.2}'s discussion of transparency sets out what transparency requires in general: documentation clear enough for peer researchers to reproduce the work and for outside bodies to audit it. Applied specifically to bias, that means publishing enough about a system's training data and evaluation methods that an outside reviewer can check for the sources of bias section~\ref{sec:6.1.1} describes, not simply take the developer's word that they were checked. That openness carries a cost — it can expose proprietary methods and hand bad actors a clearer map of how to exploit the system — a tension transparency requirements have to weigh.

\runin{COMPAS Revisited}
What COMPAS did is described above. Northpointe, the company that built it, refused to disclose the algorithm's workings, and a Wisconsin court upheld that refusal as protected trade secret even for a defendant sentenced partly on its output \autocite{loomis2016compas}. ProPublica's finding came from analyzing scores and outcomes obtained through public-records requests, not from anything the company published. A jurisdiction that required a risk-assessment tool's methodology to be audited before deployment would have caught the disparity before it touched a single defendant, not years and thousands of cases later.

\runin{Public Engagement and Oversight}
Section~\ref{sec:6.1}'s point about involving the people an AI system's decisions affect has a formal counterpart: an oversight committee with standing to monitor the system's performance, flag bias as it appears, and hold the developer accountable for fixing it — not just a channel for feedback, but a body with teeth.

\runin{Explainability Tools}
LIME and SHAP are the two tools practitioners reach for. LIME approximates a complex model's local behavior around one prediction with something simple enough to read off which inputs drove it; SHAP assigns each input a share of credit for the output, using a method borrowed from cooperative game theory \autocite{ribeiro2016why} \autocite{lundberg2017unified}. Applied to a hiring or credit-scoring model, either can surface a bias — an applicant's zip code doing more work in the decision than it should, say — that scrutinizing the training data alone would never reveal, because the bias lives in how the model combines features, not in any one feature by itself. Counterfactual explanation asks what would have had to be different: the smallest change to an applicant's financial profile that would have flipped the decision tells the applicant what mattered, and gives the institution a check on whether that factor should have mattered as much as it did. None of the three can tell you that the model is aimed at the wrong thing. All take the training target as given and report how the inputs move it, so a model whose target is itself a poor proxy for what matters — spending standing in for medical need, the case section~\ref{sec:6.3.5} works through — will attribute its predictions to the features it was built to use and report nothing amiss.
